%%%%%%%%%%%%%%%%%%%%%%%%%%%%%%%%%%%%%%%%%%%%
% FILE INI JANGAN DI UBAH ATAU MODIFIKASI
%%%%%%%%%%%%%%%%%%%%%%%%%%%%%%%%%%%%%%%%%%%%
\normalsize
\pagenumbering{roman}
\singlespacing
\hyphenpenalty=1000
\tolerance=1000
\emergencystretch=2.5em
% Cover
%\addcontentsline{toc}{chapter}{Cover}
\begin{FontCover}
	\sffamily	
\thispagestyle{empty}
\begin{figure} [!t]
	\includegraphics[keepaspectratio=true,scale=1,left]{./lib/LogoITS.png}
	\caption*{}
	\label{fig:LogoITS}
\end{figure}
\vspace{1ex}
%\noindent\makebox[\linewidth]{\rule{\paperwidth}{10pt}}
\makebox[\linewidth]{\includegraphics[keepaspectratio=true,width=\paperwidth]{./lib/\bbelt}}
\vspace{10ex}
\justify
\textbf{\Kodekk}
\vspace{1ex}
\justify
\Large
\begin{spacing}{1}
%Masukan Judul Tesis Disini
\flushleft
\textbf{\JdTesisEng}
\end{spacing}
\normalsize
\vfill
\justify

%\begin{spacing}{1}
\large
\MakeUppercase{\NamaMahasiswa}\\
NRP  \uppercase{\NrpMahasiswa}
\vspace{2ex}
\justify
\normalsize

\ifthenelse{\boolean{bMaster}}{\textbf{SUPERVISOR}}{\textbf{PROMOTOR}}\\
\PbSatu\\
\ifthenelse{\boolean{PembimbingDua}}{\ifthenelse{\boolean{bMaster}}{}{\\\textbf{Co.PROMOTOR\\}}}{}
\ifthenelse{\boolean{PembimbingDua}}{\PbDua\\}{}
\ifthenelse{\boolean{PembimbingTiga}}{\PbTiga\\}{}
 
\vfill
\justify
\bdk
\vspace{1ex}
%\rmfamily
\normalfont
\end{FontCover}
\newpage
\cleardoublepage
% Lembar pengesahan
\addcontentsline{toc}{chapter}{VALIDATION SHEET}
\ifthenelse{\boolean{bThesis}}
{ \AddToShipoutPicture*{\BackgroundIm}
	\begin{center}

		\Large\textbf{\pPengesahan}
		%\Large\textbf{PROPOSAL TESIS}
	\end{center}
	\begin{center}
		\tsss\\
		\textbf{\ggGelar}\\
		at\\
		\textbf{Institut Teknologi Sepuluh Nopember}\\
		\vspace{1ex}
		By:\\
		\textbf{\NamaMahasiswa}\\
		\textbf{NRP: \NrpMahasiswa}\\
		\vspace{1ex}
		Examination Date : \TglUjian\\
		Graduation Period : \PerWisuda\\
		\vspace{1ex}
	\end{center}
	\begin{center}
		Approved by:\\
		\textbf{Supervisor}:
	\end{center}
	\begin{enumerate}
		\item \PbSatu \hfill ...............................\\
		      NIP:\NipPbSatu \vfill
		      \ifthenelse{\boolean{PembimbingDua}}{\item \PbDua\hfill ...............................\\
		      NIP:\NipPbDua}\vfill
		      \ifthenelse{\boolean{PembimbingTiga}}{\item \PbTiga \hfill ...............................\\
		      NIP:\NipPbTiga}{}
	\end{enumerate}
	\begin{center}
		\textbf{Examiner}:
	\end{center}
	\begin{enumerate}
		\item \PjSatu \hfill ...............................\\
		      NIP:\NipPjSatu
		      \vfill
		\item \PjDua \hfill ...............................\\
		      NIP:\NipPjDua\vfill
		      \ifthenelse{\boolean{PengujiTiga}}{
		\item \PjTiga \hfill ...............................\\NIP:\NipPjTiga\vfill}{}
		      \ifthenelse{\boolean{PengujiEmpat}}{
		\item \PjEmpat \hfill ...............................\\ NIP:\NipPjEmpat\vfill}{}
	\end{enumerate}
	\vfill
	\begin{center}
		Head of \Dep \\
		\fak\\
		\vspace{9ex}
		\underline{\NmKaDep}\\
		NIP:\NipKaDep
	\end{center}
	\newpage
}
{
	\begin{center}
		\textbf{VALIDATION SHEET\\
			THESIS PROPOSAL}\\
	\end{center}
	\vspace{5ex}
	\begin{tabular}{p{2cm} c p{8cm}}
		Title & : & \JdTesisEng    \\
		By    & : & \NamaMahasiswa \\
		Nrp   & : & \NrpMahasiswa
	\end{tabular}
	\vspace{5ex}
	\begin{center}
		\textbf{Has Been Examined On}
	\end{center}
	\begin{tabular}{p{2cm} c p{8cm}}
		Day   & : & \bHariUjian   \\
		Date  & : & \TglUjian     \\
		Place & : & \bTempatUjian
	\end{tabular}
	\vspace{5ex}
	%\scalebox{.95}[1.0]{\PjSatu}

	\begin{tabular}{p{8cm} p{8cm} }
		\textbf{Examiner}                        & \textbf{Supervisor Candidate} \\
		\vspace{8ex}\hspace{-10ex}1. \PjSatu     &
		\vspace{8ex}\hspace{-8ex}1. \PbSatu                                      \\
		\hspace{-7ex}Nip :\NipPjSatu             &
		\hspace{-5ex}Nip :\NipPbSatu                                             \\
		\vspace{8ex}\hspace{-10ex}2. \PjDua      &
		\ifthenelse{\boolean{PembimbingDua}}
		{\vspace{8ex}\hspace{-8ex}2. \PbDua}{}                                   \\
		\hspace{-7ex}Nip :\NipPjDua              &
		\ifthenelse{\boolean{PembimbingDua}}
		{\hspace{-5ex}Nip :\NipPbDua}{}                                          \\
		\ifthenelse{\boolean{PengujiTiga}}
		{\vspace{8ex}\hspace{-10ex}3. \PjTiga}{} &
		\ifthenelse{\boolean{PembimbingTiga}}
		{\vspace{8ex}\hspace{-8ex}3. \PbTiga}{}                                  \\
		\ifthenelse{\boolean{PengujiTiga}}
		{\hspace{-7ex}Nip :\NipPjTiga}{}         &
		\ifthenelse{\boolean{PembimbingTiga}}
		{\hspace{-5ex}Nip :\NipPbTiga}{}                                         \\


		\ifthenelse{\boolean{PengujiEmpat}}
		{\vspace{8ex}\hspace{-10ex}4. \PjTiga}{} &                               \\
		\ifthenelse{\boolean{PengujiEmpat}}
		{\hspace{-7ex}Nip :\NipPjEmpat}{}        &                               \\
	\end{tabular}
	\newpage
}
\cleardoublepage

\ifthenelse{\boolean{bThesis}}
{
	\ifthenelse{\boolean{bMaster}}
	{
		\chapter*{STATEMENT OF THE THESIS AUTHENTICY}
		\addcontentsline{toc}{chapter}{STATEMENT OF THE THESIS AUTHENTICY}


		\begin{spacing}{1.5}

			I hereby declare that the entire content of my thesis with the title of \textbf{\JdTesisEng} is truly the result of an independent intellectual work, completed without the use of unauthorized materials, and is not the work of another party which I acknowledge as my work.

			All references cited or referred to have been written in full in the bibliography. If it turns out that this statement is not true, I am willing to accept sanctions under applicable regulations.

			\hspace{35ex}Surabaya \today

			\vspace{10ex}

			\hspace{35ex}\underline{\NamaMahasiswa}

			\hspace{35ex}Nrp: \NrpMahasiswa

		\end{spacing}
		\cleardoublepage

	}
	{

		\chapter*{PERNYATAAN KEASLIAN DISERTASI}
		\addcontentsline{toc}{chapter}{PERNYATAAN KEASLIAN DISERTASI}

		\begin{spacing}{1.5}

			Dengan ini saya menyatakan bahwa isi keseluruhan Disertasi  saya dengan judul \textbf{\JdTesis} adalah benar-benar hasil karya intelektual mandiri, diselesaikan tanpa menggunakan bahan-bahan yang tidak diijinkan dan bukan merupakan karya pihak lain yang saya akui sebagai karya sendiri.

			Semua referensi yang dikutip maupun dirujuk telah ditulis secara lengkap pada daftar pustaka. Apabila ternyata pernyataan ini tidak benar, saya bersedia menerima sanksi sesuai peraturan yang berlaku.

			\vspace{2ex}

			\hspace{35ex}Surabaya %\today

			\vspace{10ex}

			\hspace{35ex}\underline{\NamaMahasiswa}

			\hspace{35ex}Nrp: \NrpMahasiswa

		\end{spacing}
		\cleardoublepage
	}
}{}




% Kata pengantar
\newpage

\addcontentsline{toc}{chapter}{PREFACE}
\begin{center}
	\Large\textbf{PREFACE}
\end{center}
\vspace{2ex}
%Tulis kata pengantar di sini
Praise and gratitude to the presence of Allah SWT. for all His grace, the author will able to work this research with the title \textbf{Classification of Elderly Exercise Activity Based on Pose Estimation Using Deep Learning}. This research was written as a requirement for completing a master’s degree at the Department of Electrical Engineering, Institut Teknologi Sepuluh Nopember.\\
The author also expresses great thanks to:
\begin{itemize}
	\item Parents and family for their encouragement and
	      support during the author’s study.
	\item Prof. Dr. Ir. Mauridhi Hery Purnomo, M.Eng. and Dr. Eko Mulyanto Yuniarno, ST. MT. as the supervising lecturer who has provided a lot of advice on this research.
	\item Nabila Shafa Oktavia who helped the author in technical and non-technical support.
	\item Friends from the B300 laboratory who always accompany and provide support in various forms of support.
	\item As well as friends from the class of 2022 who have passed the lecture period together with the author.
\end{itemize}
\vspace{26pt}
\begin{flushright}
	\begin{tabular}[b]{c}
		Surabaya, July 2024
		\\
		\\
		\\
		The Author
	\end{tabular}
\end{flushright}

\cleardoublepage
\newpage
% Abstrak
\begin{spacing}{1}
    \begin{center}
		\large\textbf{\JdTesisEng}
	\end{center}
	\normalsize
	\begin{adjustwidth}{-0.2cm}{}
		\begin{tabular}{lcp{0.9\linewidth}}
		By &:& \NamaMahasiswa\\
			Student Identity Number &:&\NrpMahasiswa\\
			Supervisor &:& 1. \PbSatu\\
			\ifthenelse{\boolean{PembimbingDua}}{& & 2. \PbDua\\}{}
			\ifthenelse{\boolean{PembimbingTiga}}{& & 3. \PbTiga\\}{}
		\end{tabular}
	\end{adjustwidth}
	\vspace{2ex}
	
	\begin{center}
		\Large\textbf{ABSTRACT}
	\end{center}
	\vspace{1ex}
    %Tulis Abstrak disini
    Older people's productivity often declines, especially in terms of physical ability. Physical decline can be slowed down with exercise and other physical activities. However, activities such as stretching are often neglected by the elderly. Exercise activities for the elderly are important so that the elderly can maintain their health in old age. Research on elderly activity recognition has been widely developed. Convolutional Neural Network (CNN) is a type of artificial neural network specifically designed for image processing and recognition. On the other hand, Long Time-Term Memory (LSTM) is an efficient method for solving real-time problems. These two methods can be used for labeling and recognizing physical activity movements in the elderly. Physical activities performed by the elderly need to be customized. In this study, we developed a model with elderly pose estimation. One of the frameworks for human pose estimation is Mediapipe Pose Estimation (MPE). Therefore, this research focuses on the recognition and detection of fitness movements in the elderly. The work begins with the acquisition of datasets in the form of exercise activities that have been adapted to the physical issues of the elderly. Data acquisition was carried out since there are few and limited datasets that discuss this training activity. The acquired video data was then subjected to a video frame extraction process. Each frame sequence represents the exercise activity information. Pose estimation has been done using the Mediapipe framework. The extraction results are then trained using CNN, LSTM, CNN-LSTM, and deep CNN-LSTM architectures. The accuracy of each model is 83.68\%, 92.89\%, 96.05\%, and 87.11\%. Based on these results, the CNN-LSTM model outperforms the other models with an accuracy rate of 96.05\%. The error in recognizing data patterns is shown using the loss metric. The loss value of the CNN-LSTM model is 0.1498, the smallest compared to other models. This value indicates the model's ability to predict data with the lowest error rate. In addition, in other metrics, this model outperforms other models. Precision, recall, and f1-score of the CNN-LSTM model are at 0.96, respectively.

    %Tulis Kata Kunci disini
    \vspace{2ex}
    \textbf{Keyword}: Activity, Deep Learning, Elderly, Pose Estimation
\end{spacing}
\cleardoublepage
\ifthenelse{\boolean{JudulTesisEng}}
{
	\begin{spacing}{1}
    \begin{center}
		\large\textbf{\JdTesis}
	\end{center}
	\normalsize
	\begin{adjustwidth}{-0.2cm}{}
		\ifthenelse{\boolean{bMaster}}{
		
		\begin{tabular}{lcp{0.9\linewidth}}
		Nama Mahasiswa &:& \NamaMahasiswa\\
			NRP &:&\NrpMahasiswa\\
			Pembimbing &:& 1. \PbSatu\\
			\ifthenelse{\boolean{PembimbingDua}}{& & 2. \PbDua\\}{}
			
			\ifthenelse{\boolean{PembimbingTiga}}{& & 3. \PbTiga\\}{}
			
			
			
		\end{tabular}
	}{
	\begin{tabular}{lcp{0.7\linewidth}}
		Nama Mahasiswa &:& \NamaMahasiswa\\
		NRP &:&\NrpMahasiswa\\
		Promotor &:&  \PbSatu\\
		\ifthenelse{\boolean{PembimbingTiga}}{
		\ifthenelse{\boolean{PembimbingDua}}{Co. Promotor&: & 1. \PbDua\\}{}
		\ifthenelse{\boolean{PembimbingTiga}}{& & 2. \PbTiga\\}{}
	}
{
	\ifthenelse{\boolean{PembimbingDua}}{\hspace{5ex}Co. Promotor&: & \PbDua\\}{}
	
}
	\end{tabular}

}

	
	\end{adjustwidth}
	\vspace{2ex}
	\begin{center}
		\Large\textbf{ABSTRAK}
	\end{center}
	\vspace{1ex}
    %Tulis Abstrak disini
    Produktivitas orang tua sering kali menurun, terutama dalam hal kemampuan fisik. Penurunan kemampuan fisik dapat diperlambat dengan olahraga dan aktivitas fisik lainnya. Namun, aktivitas seperti peregangan sering diabaikan oleh orang tua. Aktifitas latihan untuk elderly ini menjadi penting agar elderly dapat menjaga kesehatannya di usia lanjut. Riset mengenai pengenalan aktivitas lansia telah banyak dikembangkan. Convolutional Neural Network (CNN) adalah jenis jaringan saraf tiruan yang dirancang khusus untuk pengolahan dan pengenalan gambar. Di sisi lain, Long Shoert-Term Memory (LSTM) adalah metode yang efisien untuk memecahkan masalah real-time. Kedua metode ini dapat digunakan untuk pelabelan dan pengenalan gerakan aktivitas fisik pada lansia. Aktivitas fisik yang dilakukan oleh lansia perlu disesuaikan. Dalam penelitian ini, kami mengembangkan sebuah model dengan estimasi pose lansia. Salah satu kerangka kerja untuk estimasi pose manusia adalah Mediapipe Pose Estimation (MPE). Oleh karena itu, penelitian ini berfokus pada pengenalan dan deteksi gerakan kebugaran pada lansia. Pekerjaan diawali dengan akuisisi dataset berupa aktifitas latihan yang telah disesuaikan dengan isu-isu fisik para elderly. Akusisi data dilakukan sejak sedikit dan terbatasnya dataset yang membahas aktifitas latihan ini. Data video yang telah diakuisisi kemudian dilakukan proses ekstraksi video frame. Setiap urutan frame mewakili informasi aktifitas latihan. Estimasi pose telah dilakukan menggunakan framework Mediapipe. Hasil ekstraksi ini kemudian dilatih menggunakan arsitektur CNN, LSTM, CNN-LSTM, dan deep CNN-LSTM. Akurasi setiap model sebesar 83.68\%, 92.89\%, 96.05\%, dan 87.11\%. Berdasarkan hasil tersebut, model CNN-LSTM mengungguli model-model lainnya dengan tingkat akurasi 96.05\%. Kesalahan dalam mengenali pola data ditunjukkan menggunakan metric loss. Nilai loss model CNN-LSTM sebesar 0.1498, paling kecil dibandingkan dengan model-model lainnya. Nilai ini menunjukkan kemampuan model dalam memprediksi data dengan tingkat kesalahan paling rendah. Selain itu, pada metrcis lainnya, model ini mengungguli daripada model lainnya. Precision, recall, dan f1-score model CNN-LSTM berada pada nilai 0.96, masing-masing.

    %Tulis Kata Kunci disini
    \vspace{2ex}
    \textbf{Kata kunci }: Aktivitas, \textit{Deep Learning}, Lansia, Estimasi Pose

\end{spacing}
}{}
\cleardoublepage

% Daftar isi
\renewcommand*\contentsname{TABLE OF CONTENTS}
\titlespacing*{\chapter}{0pt}{0ex}{0ex}
\tableofcontents
\cleardoublepage
% Daftar gambar
\renewcommand*\listfigurename{LIST OF FIGURES}
\addcontentsline{toc}{chapter}{\listfigurename}
\titlespacing*{\chapter}{0pt}{0ex}{0ex}
\listoffigures
\cleardoublepage
% Daftar tabel
\renewcommand*\listtablename{LIST OF TABLES}
\addcontentsline{toc}{chapter}{\listtablename}
\titlespacing*{\chapter}{0pt}{0ex}{0ex}
\listoftables
\cleardoublepage
% Nomenklatur
\addcontentsline{toc}{chapter}{NOMENCLATURE} % kata pengantar
\begin{center}
	\Large\textbf{NOMENCLATURE}
\end{center}
\vspace{1ex}
\begin{tabular}{ll}
	$b$       & LSTM - bias                         \\
	Conv1D    & Convolutional 1 Dimension           \\
	CNN       & Convolution Neural Network          \\
	$C_t$     & LSTM - Cell State                   \\
	FC        & Fully Connected                     \\
	FN        & False Negative                      \\
	FP        & False Positive                      \\
	$f_t$     & LSTM - Forget Gate                  \\
	$h_t$     & LSTM - Hidden Outputs               \\
	$h_{t-1}$ & LSTM - Hidden State at $t$ - 1 time \\
	$i_t$     & LSTM - Input Gate                   \\
	LSTM      & Long Short Term Memory              \\
	NN        & Neural Networks                     \\
	$o_t$     & LSTM - Output Gate                  \\
	ReLU      & Rectified Linear Unit               \\
	RNN       & Recurrent Neural Networks           \\
	$\sigma$  & LSTM - Sigmoid Function             \\
	tanh      & LSTM - tanh Function                \\
	TN        & True Negative                       \\
	TP        & True Positive                       \\
	$V$       & Mediapipe - Keypoint Vector         \\
	W         & Weights                             \\
	$x_t$     & LSTM - Input at t time              \\
	$z_{max}$ & Mediapipe - Maximum Depth           \\
	$z_{min}$ & Mediapipe - Minimum Depth           \\

	% $s_i$            & Static Variables                                \\
	% $s^{2}$          & Sample Variance                                 \\
	% TFT              & Temporal Fusion Transformers                    \\
	% VSN              & Variable Selection Network                      \\
	% $\boldsymbol{v}$ & Attention - Value    							 \\
	% $W_H$            & Head Weight                                     \\
	% $W_K$            & Key Weight                                      \\
	% $W_Q$            & Query Weight                                    \\
	% $W_V$            & Value Weight                                    \\
\end{tabular}
\newpage

\cleardoublepage

% Bab isi buku
\titleformat{\chapter}[display]{\bfseries\Large}{CHAPTER \centering\thechapter}{0ex}{\vspace{0ex}\centering}[\vspace{0ex}]
\titleformat{\section}{\bfseries\large}{\MakeUppercase{\thesection}}{1ex}{}
\titleformat{\subsection}{\bfseries\large}{\MakeUppercase{\thesubsection}}{1ex}{}
\titleformat{\subsubsection}{\bfseries\large}{\MakeUppercase{\thesubsubsection}}{1ex}{}
\titlespacing*{\chapter}{0pt}{-4ex}{0pt}
\titlespacing{\section}{0pt}{0pt}{0pt}
\titlespacing{\subsection}{0pt}{0pt}{0pt}
\titlespacing{\subsubsection}{0pt}{0pt}{0pt}
% Indent paragraph
\setlength{\parindent}{0.8cm}